\section{Privacy purpose and the lifecycle of memory}\label{sec:C29:privacy}
Privacy depends on collection, use, sharing, retention and deletion, not only on whether storage is encrypted. An agent may retrieve a permitted document for one purpose and later reuse its summary in another context where that use is inappropriate. A memory record therefore needs subject, source, purpose, audience, validity and retention metadata in addition to content. These fields influence both retrieval and action authorization.

Purpose limitation is difficult to infer from natural language alone. A broad request to help with a project does not necessarily authorize exporting every related document. Application logic should define the available capability envelope, and the agent should request additional authority only when the task requires it. The request should identify the specific effect and evidence, not seek a blanket expansion that persists indefinitely.

Deletion has several meanings. Logical exclusion prevents future retrieval; physical removal deletes stored bytes; derivative invalidation addresses summaries, embeddings and skills; parameter unlearning concerns information absorbed into model weights. Completing one does not establish the others. A system should accurately report which action occurred and which copies or retention obligations remain. External disclosure cannot generally be undone by deleting the local source.

Shared memory raises multi-user conflicts. A procedure learned from one user's private task may be useful to another, but utility does not grant permission to share it. Abstraction can still reveal sensitive details or encode privileged assumptions. A promotion process should check whether the generalized skill contains source-specific information and whether its reuse is authorized. Evaluation should include cross-user leakage and benign transfer rather than reward sharing alone.

Privacy-aware traces require a similar balance. Detailed records support auditing and reproducibility but can expose the very content the system should protect. Separate access-controlled raw evidence from redacted operational metadata, define retention, and preserve stable identities so authorized auditors can join records. If redaction removes decisive evidence from public artifacts, state the resulting verification limit rather than imply full independent reproducibility.

